\section{Packaging: single-file executables for Linux, Windows and macOS}
\label{sec:packaging}

\subsection{Principles}

The code is pure Python with no compiled components of its own, so
``porting'' is a packaging exercise, done with PyInstaller. Three facts
shape everything below. First, \textbf{PyInstaller does not
cross-compile}: a Windows \code{.exe} must be built on Windows, a macOS
app on macOS, a Linux binary on Linux --- one build per platform, all
from the same source and the same spec. Second, in \code{--onefile} mode
the executable unpacks itself to a temporary directory at every start;
bundled data is found under \code{sys.\_MEIPASS}, \emph{not} next to the
executable, so path handling needs the small shim below. Third, the
program \emph{writes} beside itself (\code{etc\_user\_config.json},
\code{etc\_sites.json}, imported templates) --- those files must live in
a real, persistent, user-writable directory, never inside the bundle.

The clean policy, and the one assumed here: \textbf{bundle the read-only
\code{data/} directory inside the executable, keep the writable files
beside the executable}. Add this resolver near the top of
\code{etc\_gui.py} (and use it wherever the default data directory is
composed):

\begin{verbatim}
import sys
from pathlib import Path

def bundled_data_dir():
    """data/ inside a PyInstaller bundle, or beside the source."""
    if getattr(sys, "frozen", False):        # running from PyInstaller
        return Path(sys._MEIPASS) / "data"
    return Path(__file__).resolve().parent / "data"
\end{verbatim}

The GUI's data-directory field keeps working as the override for users
with their own data trees.

\subsection{Common build steps}

On each platform, in a fresh virtual environment:

\begin{verbatim}
python3 -m venv build-env
. build-env/bin/activate        # Windows: build-env\Scripts\activate
pip install -r requirements.txt pyinstaller
pyinstaller --onefile --windowed --name spetc \
    --add-data "data:data" \
    --add-data "etc_sites.json:." \
    --add-data "observatories.json:." \
    etc_gui.py
\end{verbatim}

(On Windows the \code{--add-data} separator is a semicolon:
\code{"data;data"}.) The result lands in \code{dist/}. PyInstaller's
built-in hooks handle NumPy, SciPy, Matplotlib and Astropy; two points
need attention. Astropy must not try to download IERS tables from a
frozen app --- the code already sets \code{iers.conf.auto\_download =
False}, so the bundled table is used; verify by running the frozen app
with the network off. And Matplotlib should be pinned to the TkAgg
backend in the frozen app (\code{MPLBACKEND=TkAgg} via the spec's
\code{runtime\_hooks} or an environment default) to avoid backend probing
at start-up.

Always test the frozen build with the checklist: starts offline; template
list populated (bundled data found); a photometry and a spectroscopy run
complete; \textsf{Save instrument profile} and site saving write real
files next to the executable; the two CSV exports open.

\subsection{Linux}

The onefile binary from the recipe above runs on any distribution with
glibc at least as new as the build machine's --- so \textbf{build on the
oldest distribution you want to support} (a container of Debian
oldstable or Ubuntu LTS is the usual choice). Tk is bundled by
PyInstaller. Users make it executable (\code{chmod +x spetc}) and run
it. For desktop integration ship a \code{.desktop} file, or go one step
further and wrap \code{dist/spetc} into an AppImage with
\code{appimagetool}, which adds the icon/menu integration without
changing anything in the build.

\subsection{Windows}

Build with the python.org interpreter (not the Microsoft Store one,
whose sandbox breaks file dialogs and per-user paths). The recipe yields
\code{dist\textbackslash spetc.exe}; \code{--windowed} suppresses the
console window. Two Windows realities to plan for: SmartScreen will warn
on an unsigned executable --- users click ``More info $\to$ Run
anyway'', or you sign with an Authenticode certificate; and some
antivirus products false-positive on PyInstaller onefile stubs --- if
reports come in, prefer \code{--onedir} (a folder instead of a single
file, same command without \code{--onefile}, zipped for distribution),
which is far less often flagged. High-DPI rendering is handled by Tk on
Python $\geq$ 3.12; on older versions call
\code{ctypes.windll.shcore.SetProcessDpiAwareness(1)} at start-up if the
UI looks blurry.

\subsection{macOS}

Build on the oldest macOS you want to support, with the python.org
universal2 interpreter (bundles its own Tcl/Tk; never use the system
Python). With \code{--windowed} PyInstaller produces both a UNIX binary
and a \code{spetc.app} bundle --- distribute the \code{.app}, zipped.
Gatekeeper will quarantine an unsigned download: users right-click
$\to$ \emph{Open} once, or remove the attribute
(\code{xattr -dr com.apple.quarantine spetc.app}); for friction-free
distribution you need an Apple Developer ID, \code{codesign
--deep} and notarization with \code{xcrun notarytool} --- mechanical
but paid. An app built on Apple Silicon runs on Intel Macs only through
Rosetta if built \code{universal2}; check
\code{lipo -archs dist/spetc.app/Contents/MacOS/spetc} if you support
both.

\subsection{Reproducibility}

Commit the generated \code{spetc.spec} once it works --- it \emph{is}
the build configuration --- and record in the release notes the Python
and PyInstaller versions of each artefact. A small
\code{build.sh}/\code{build.bat} that creates the venv, installs pinned
versions and runs PyInstaller turns the release build into one command
per platform.
